\documentclass[12pt]{article}
\usepackage[utf8]{inputenc}
\usepackage[T1]{fontenc}
\usepackage{hyperref}
\title{Yangho Language Resources for K-12 Curriculum Development}
\author{Compiled by Claude Sonnet 5.5 (claude-sonnet-5-5) for Andr\'as Kornai}
\date{}
\begin{document}
\maketitle

\section{Introduction}
Yangho (ISO 639-3 ynh; Glottocode yang1292) is, according to Glottolog, a code that was retired in 2015 as non-existent. It was listed for Gabon, around Mamidi and Bakoumba. Glottolog reports that linguists working in Gabon do not recognise such a language and place Njebi (nzb) at Bakoumba instead. A German Wikipedia article still describes it as a Bantu language with about 5,000 speakers in Gabon (1990 census), under the alternate names Bayongho, Miyangho and Yongho. These sources conflict, and the existence of a distinct Yangho language is doubtful. It should not be confused with Yangben (yav) of Cameroon.

\section{Primary Educational Resources}
None found.

\section{Linguistic and Reference Materials}
\begin{itemize}
\item Glottolog entry for Yangho: \url{https://glottolog.org/resource/languoid/id/yang1292}
\item Wikipedia article: \url{https://en.wikipedia.org/wiki/Yangho}
\end{itemize}

\section{Religious and Community Resources}
None found.

\section{Digital and Audio Resources}
None found.

\section{Curriculum Development Context}
There is no basis in the search results for a Yangho-specific curriculum. If speakers in the Bakoumba area are in fact Njebi speakers, any work should start from Njebi (nzb) resources, which I did not search for here.
\end{document}